\section{代表性攻击路径二：自动化攻击生成与多 Agent 利用}

讲者展示了自动化红队思路：让 analyzer agent 先分析目标代理代码和工具接口，再由另一个 agent 自动生成攻击候选（例如上下文感知注入、命令诱导注入）。这种 ``agent 攻 agent'' 模式说明攻击者也会利用 AI 自动化其攻击开发流程。

\begin{knowledgebox}{为什么自动化攻击值得重视}
自动化攻击将显著降低攻击门槛：
\begin{itemize}
    \item 更快发现注入点和权限路径；
    \item 更容易进行批量变体生成与回归测试；
    \item 更容易绕过静态规则，因为攻击可按上下文动态调整。
\end{itemize}
\end{knowledgebox}

课程中提到的 GitHub issue solving agent 场景尤其典型：攻击者将恶意指令嵌入 issue 内容，Agent 在自动处理时执行了不应执行的命令，导致攻击链成立。这个案例突出了两个关键问题：其一，``看起来像任务描述'' 的文本很难靠关键词拦截；其二，工具执行阶段如果没有强约束，计划偏移会直接转化为高危动作。

\begin{warningbox}{白盒假设下的防御压力}
一旦攻击者可见提示词、策略逻辑或工具配置，其攻击成功率会显著提升。此时仅靠 obscurity 不可持续，必须依赖可验证策略和运行时 enforcement。
\end{warningbox}

\subsection{本章小结}
自动化攻击框架表明：防御者必须假设攻击者也在使用 Agent。安全评测不能只测单次样例，而要测连续适应性对抗。

